\section{Multi-GPU Computing with Python} 

Many Python applications can use multiple GPUs simultaneously. 

Typical examples include: 

\begin{itemize} 
\item PyTorch 
\item TensorFlow 
\item JAX 
\item MPI-based scientific computing 
\end{itemize} 

\subsection{Loading Required Modules} 

Load the required software: 

\begin{lstlisting} 
module load cuda/12.3.2 
module load anaconda3/2024.02.1 
module load openmpi/4.1.7-with-cuda 
\end{lstlisting} 

Verify: 

\begin{lstlisting} 
module list 
\end{lstlisting}

\subsection{Installing Required Packages} 

If necessary: 

\begin{lstlisting} 
pip install mpi4py 
pip install torch 
\end{lstlisting}

\subsection{Testing CUDA} 

Create a file called \texttt{test-cuda.py}: 

\begin{lstlisting} 
import torch 
print(torch.cuda.is_available()) 
print(torch.cuda.device_count()) 
\end{lstlisting}

Run: 

\begin{lstlisting} 
python test_cuda.py 
\end{lstlisting} 

Expected output: 
\begin{lstlisting} 
True 
4 
\end{lstlisting}

\subsection{Verifying GPU Assignment} 

Create a file called \texttt{test-multigpu.py}: 

\begin{lstlisting} 
from mpi4py import MPI 
import torch 
comm = MPI.COMM_WORLD 
rank = comm.rank 
torch.cuda.set_device(rank) 
print( f"Rank {rank} " f"using GPU " f"
{torch.cuda.current_device()}" ) 
\end{lstlisting}

Run four MPI ranks: 

\begin{lstlisting} 
mpiexec -n 4 python test-multigpu.py 
\end{lstlisting}

\subsection{Selecting GPUs} 

A program can be restricted to specific GPUs using \texttt{CUDA\_VISIBLE\_DEVICES}. 

Use only GPU 0: 

\begin{lstlisting} 
CUDA_VISIBLE_DEVICES=0 python script.py 
\end{lstlisting} 

Use GPUs 1 and 2: 

\begin{lstlisting} 
CUDA_VISIBLE_DEVICES=1,2 python script.py 
\end{lstlisting} 

This is useful on shared servers where multiple users may be using GPUs simultaneously. 

\subsection{MPI with Python} 

Many scientific Python applications use MPI through \texttt{mpi4py}. 

Install: 

\begin{lstlisting} 
pip install mpi4py 
\end{lstlisting} 

Create a file called \texttt{test-mpi.py}: 

\begin{lstlisting} 
from mpi4py import MPI 

comm = MPI.COMM_WORLD 

print( f"Rank {comm.rank} " f"of {comm.size}" ) 
\end{lstlisting} 

Run: 

\begin{lstlisting} 
mpiexec -n 4 python test_mpi.py 
\end{lstlisting} 

Expected output: 

\begin{lstlisting} 
Rank 0 of 4 
Rank 1 of 4 
Rank 2 of 4 
Rank 3 of 4 
\end{lstlisting} 

\subsection{Multi-GPU PyTorch} 

PyTorch supports multiple GPUs using Distributed Data Parallel (DDP). 

Check visible GPUs: 

\begin{lstlisting} 
import torch 

print(torch.cuda.device_count()) 
\end{lstlisting} 

A typical launch command is: 

\begin{lstlisting} 
torchrun --nproc_per_node=4 train.py 
\end{lstlisting} 

where one process is launched for each GPU. 

\subsection{Troubleshooting} 

Verify that Python can see the GPUs: 

\begin{lstlisting} 
nvidia-smi 
\end{lstlisting} 

Verify that PyTorch detects CUDA: 

\begin{lstlisting} 
python -c "import torch; print(torch.cuda.is_available())" 
\end{lstlisting} 

Expected output: 

\begin{lstlisting} 
True 
\end{lstlisting}